\documentclass[11pt,a4paper]{article}
\usepackage[margin=2cm]{geometry}
\usepackage[T1]{fontenc}
\usepackage{lmodern}
\usepackage{amsmath}
\usepackage{booktabs}
\usepackage{tabularx}
\usepackage{graphicx}
\usepackage{enumitem}
\usepackage{xcolor}
\usepackage[colorlinks=true,linkcolor=blue!50!black,urlcolor=blue!50!black]{hyperref}
\setlist{nosep,leftmargin=*}
\setlength{\parskip}{4pt}
\setlength{\parindent}{0pt}
\newcommand{\tbd}[1]{\textcolor{red}{[TBD: #1]}}
\newcommand{\owner}[1]{\textcolor{gray}{\small(#1)}}
\newcommand{\figfile}[2]{\IfFileExists{../#1}{\includegraphics[width=#2]{../#1}}{\IfFileExists{#1}{\includegraphics[width=#2]{#1}}{\fbox{\parbox[c][3cm][c]{#2}{\centering\tbd{figure \texttt{\detokenize{#1}}}}}}}}
\newcommand{\teamcontacts}{}
\IfFileExists{team_private.tex}{\input{team_private.tex}}{\IfFileExists{docs/team_private.tex}{\input{docs/team_private.tex}}{}}

\title{\textbf{Industrial Defect Segmentation on MVTec AD}\\[2pt]\large Stage 2: First Working Model -- Computer Vision 2026}
\author{Team: jestersw (lead), moonshine1l, psandersi, AigulFar}
\date{13 October 2026\\[2pt]\small Repository: \url{https://github.com/jestersw/Defect-Segmentation} \quad Commit: \tbd{hash}}

\begin{document}
\maketitle
\vspace{-1.5em}

\section{Summary} \owner{jestersw}

We implemented the data loading pipeline, a U-Net baseline trained with binary cross-entropy (BCE) and the evaluation code, and trained the baseline on the 550-image training split at $256\times256$ on CPU. On the validation split the baseline reaches Dice \tbd{x.xx} and small-defect recall \tbd{x.xx}. The test split has not been used; it is reserved for the final comparison with the Dice + Focal condition.

\section{Data Loading and Preprocessing} \owner{moonshine1l}

\textbf{Dataset class.} \texttt{MVTecSegDataset} (\path{src/data.py}) reads the fixed manifests in \path{splits/} (550 / 118 / 118 images). For every image it loads the ground-truth mask from \path{<category>/ground_truth/<defect_type>/<stem>_mask.png}; good images receive an all-zero mask. Each item is an image tensor $[3,256,256]$ and a binary mask tensor; masks are $[1,256,256]$ for training and $[1,512,512]$ for validation and test, so that metrics are computed at the resolution of the small-defect threshold (859.28 px). Both tensors are contiguous \texttt{float32}; \texttt{records} preserves manifest order for category-wise evaluation.

\textbf{Preprocessing.} Images are read as RGB, resized with bilinear interpolation and normalised with ImageNet mean and standard deviation. Masks are resized directly from their original resolution with nearest-neighbour interpolation and binarised at $>0$; evaluation masks never pass through the smaller input resolution.

\textbf{Augmentation (train only).} Horizontal and vertical flips and $90^\circ$ rotations are applied jointly to image and mask. Image-only \texttt{RandomBrightnessContrast} uses brightness and contrast limits of 0.1; the brightness offset is relative to the maximum pixel value. Each of the four transforms has probability 0.5. Validation and test are deterministic.

\textbf{Loaders and tests.} Batch size 8, four workers, and only training data shuffled; the last partial batch is retained. Independent seeded generators and worker-specific augmentation seeds reproduce runs under the same worker count and library versions. All 26 data tests pass on synthetic images, covering tensor formats, RGB normalisation, binary/empty masks, paired augmentation, deterministic evaluation, manifest order, Windows-compatible worker spawning and integration with evaluation at a different mask resolution. A local check read all 550 training images in 69 batches and 118 validation images in 15 batches; no real test images were loaded. One loader pass took 21.87 s for train and 6.54 s for validation, including worker startup and tensor checks, without a model. These are local diagnostic timings under Python 3.13 / PyTorch 2.10, rather than the project Python 3.11 / PyTorch 2.6 baseline environment. Settings, versions and split hashes are recorded in \path{results/s2/data_loading_check.json}.

\begin{figure}[htbp]
\centering
\figfile{results/s2/data_samples.png}{0.85\linewidth}
\caption{Training samples after augmentation with ground-truth masks overlaid.}
\label{fig:data}
\end{figure}

\section{Baseline Model and Training Setup} \owner{psandersi}

\begin{table}[h]
\centering
\small
\begin{tabularx}{\linewidth}{lX}
\toprule
Model & U-Net (\texttt{segmentation\_models\_pytorch} 0.4.0), ResNet34 encoder with ImageNet weights, 1 output channel (logits) \\
Parameters & \tbd{number} \\
Loss & \texttt{BCEWithLogitsLoss} \\
Optimiser & AdamW, learning rate $10^{-4}$, weight decay $10^{-4}$, cosine schedule \\
Input / batch & $256\times256$, batch size 8, 69 steps per epoch \\
Epochs & at most 60, early stopping on validation Dice with patience 10 \\
Seed & 42, deterministic algorithms \\
Hardware & laptop CPU (AMD Ryzen AI 9 HX 370, 31\,GiB RAM), PyTorch 2.6.0 \\
Training time & \tbd{min per epoch, total time, best epoch} \\
\bottomrule
\end{tabularx}
\caption{Baseline configuration (\texttt{results/runs/baseline\_bce\_256/config.yaml}).}
\end{table}

\textbf{Training procedure.} \path{src/train.py} trains the model from \path{configs/base.yaml}, evaluates on the validation split after every epoch and keeps the checkpoint with the highest validation Dice. Per-epoch losses, validation metrics and timings are written to \path{results/runs/baseline_bce_256/history.csv}.

\begin{figure}[htbp]
\centering
\figfile{results/runs/baseline_bce_256/curves.png}{0.8\linewidth}
\caption{Training and validation loss and validation Dice per epoch.}
\label{fig:curves}
\end{figure}

\section{Initial Results (Validation Split)} \owner{jestersw}

\begin{table}[h]
\centering
\small
\begin{tabular}{lrrrrrr}
\toprule
 & Dice & IoU & Pixel P & Pixel R & Small-defect R & Good FPR \\
\midrule
All      & \tbd{} & \tbd{} & \tbd{} & \tbd{} & \tbd{} & \tbd{} \\
\midrule
capsule  & \tbd{} & \tbd{} & \tbd{} & \tbd{} & \tbd{} & \tbd{} \\
hazelnut & \tbd{} & \tbd{} & \tbd{} & \tbd{} & \tbd{} & \tbd{} \\
metal\_nut & \tbd{} & \tbd{} & \tbd{} & \tbd{} & \tbd{} & \tbd{} \\
tile     & \tbd{} & \tbd{} & \tbd{} & \tbd{} & \tbd{} & \tbd{} \\
wood     & \tbd{} & \tbd{} & \tbd{} & \tbd{} & \tbd{} & \tbd{} \\
\bottomrule
\end{tabular}
\caption{BCE baseline on the validation split (118 images), threshold 0.5, predictions upsampled to $512\times512$. Source: \texttt{results/runs/baseline\_bce\_256/metrics\_val.json}.}
\label{tab:results}
\end{table}

\textbf{Threshold sensitivity.} At thresholds 0.3 / 0.5 / 0.7 the pixel recall is \tbd{} / \tbd{} / \tbd{} and the pixel precision \tbd{} / \tbd{} / \tbd{}.

\textbf{Interpretation.} \tbd{2--3 sentences: which categories are easy or hard, how small-defect recall compares with overall recall, whether the model predicts defects on good images}

\begin{figure}[htbp]
\centering
\figfile{results/s2/baseline_examples.png}{0.9\linewidth}
\caption{Validation examples: image, ground truth and baseline prediction, including small defects and one good image.}
\label{fig:examples}
\end{figure}

\section{Problems and Observations} \owner{everyone}

\begin{itemize}
  \item \textbf{CPU-only compute.} No GPU is available, so training runs at $256\times256$ with one seed; a 60-epoch run at $512\times512$ would take about 9\,h.
  \item \textbf{Validation only.} Initial results use the validation split; the test split is kept for the final comparison to avoid selection bias.
  \item \textbf{Evaluation mask resolution.} A synthetic one-pixel defect survives in the evaluation mask even when the model input is smaller: masks are resized independently from the original annotation. This keeps ground truth at the resolution of the fixed small-defect threshold. The lower-resolution model may still miss such defects.
  \item \tbd{psandersi: training observation (convergence, time per epoch vs.\ smoke test, memory)}
  \item \tbd{jestersw: results observation (e.g.\ missed small defects, false positives)}
  \item \tbd{AigulFar: observation on the losses}
\end{itemize}

\section{Next Steps} \owner{AigulFar}

\begin{itemize}
  \item Train condition B (Dice + Focal) with the same configuration and compare on validation (stage 3, 20 Oct).
  \item Retrain the better condition at $512\times512$ for the resolution ablation.
  \item \tbd{anything learned from the baseline that changes the plan}
\end{itemize}

\textbf{Contributions.} jestersw: evaluation, results, report; moonshine1l: data pipeline; psandersi: model, training loop, baseline run; AigulFar: loss functions.

\teamcontacts

\end{document}
